%% kapa: exam=1 round=35 year=2024 period=1 form=A subject=real_estate subject_ko=부동산학원론 seq=3 number=83
%% source: Q-Net 감정평가사 기출문제 (cst00309) artlSeq=5213229 file=raw/1008/5213229/1교시 문제지 원본.pdf
%% conversion: kordoc 4.24.51 -> tools/gen1.py (auto); review=proofread-sonnet; flags=none
\begin{question}{83}
주택법령상 주택의 정의에 관한 설명으로 옳은 것은?
\begin{choices}
  \item 민영주택은 임대주택을 제외한 주택을 말한다.
  \item 세대구분형 공동주택은 공동주택의 주택 내부 공간의 일부를 세대별로 구분하여 생활이 가능한 구조로 하되, 그 구분된 공간의 일부를 구분소유 할 수 있는 주택으로서 대통령령으로 정하는 건설기준, 설치기준, 면적기준 등에 적합한 주택을 말한다.
  \item 도시형 생활주택은 300세대 미만의 국민주택규모에 해당하는 주택으로서 대통령령으로 정하는 주택을 말한다.
  \item 에너지절약형 친환경주택은 저에너지 건물 조성기술 등 대통령령으로 정하는 기술을 이용하여 에너지 사용량을 절감하거나 이산화탄소 배출량을 증대할 수 있도록 건설된 주택을 말한다.
  \item 장수명 주택은 구조적으로 오랫동안 유지ㆍ관리될 수 있는 내구성을 갖추고 있어 내부 구조를 쉽게 변경할 수 없는 주택을 말한다.
\end{choices}
\end{question}
\answer{3}
